\section{Discussion}
\label{sec:discussion}

\subsection{What the Price of the Unit Depends On}
The central finding (RQ1, RQ2) is not that an individual-level guarantee has a cost,
which was expected, but that the cost has a regular shape. At one operating
point on person-level VitalDB it is $0.0170$ AUROC. Holding the per-step budget and the units per step
fixed, so that only the clipping geometry differs, it is indistinguishable
from zero at one record per individual and grows as $A(1-1/k)$, with a single
parameter that a fit on $k \le 4$ already determines. The patient-level model
stays flat while the record-level model improves, so a patient-level guarantee
becomes relatively more expensive, not cheaper, as more records are collected
per patient. Normalizing the clipping bound by the record count does not
recover the cost on eICU and could explain at most one eighth of it on VitalDB,
so the cost belongs mainly to the unit. Proposition~\ref{prop:scaling}
also points to what would: clipping around an estimate of each individual's
mean gradient, whose user-level rates improve with $k$ \cite{ref32}. The
external corpus reproduces the shape and the saturation but not the
magnitude: the relationship transfers in shape, its magnitude is
corpus-dependent, and outcome prevalence, client heterogeneity, contribution
counts and feature distributions are plausible contributors that the present
experiments do not isolate.

\subsection{Position Against Prior Art}
PP-HFM builds on user-level privacy \cite{ref11}
and parameter-efficient private fine-tuning \cite{ref12, ref13}; its
contribution is their composition inside a hierarchy, with a proof, and its
cost. Hierarchical systems
\cite{ref3, ref4, ref5} report savings without a guarantee, and private
systems \cite{ref9, ref10} leave the edge tier out of the accounting; PP-HFM
measures both on the same configuration and shows that the savings cost no
utility. The re-implemented baselines make this concrete
(RQ4; Section~\ref{sec:results:prior}). On eICU, HED-FL matches the non-private
reference, so its $3.2$-point lead over PP-HFM is the price of the
patient-level guarantee; on VitalDB, PP-HFM matches HED-FL in AUROC, exceeds
it in AUPRC, and leaks less under offline LiRA (margin $-0.005$ against $+0.026$). FL2DP adds noise yet leaks more than HED-FL, which adds none:
a perturbation without a finite $\varepsilon$ is not a weaker guarantee but
no guarantee, and only an audit can show what it protects. The parameter-efficient literature attributes
robustness at small $\varepsilon$ to the reduced privatized dimension; the
three-way decomposition attributes it to the initialization instead (RQ3), which
changes what a practitioner should optimize first. Several results are
negative and are reported with the control or the power that makes them
interpretable: the privatized subset produces no resolved accuracy effect once the
initialization is controlled, compression and loss-aware weighting produce no
resolved accuracy change, and count-normalized clipping does not recover the
cost of the unit on eICU. Within the evaluated configurations, they indicate
where engineering effort does not pay off in accuracy.

\subsection{Practical Implications}
Four recommendations follow. First, choose the unit explicitly and price it
in records per individual: where individuals contribute many records, a
record-level budget bounds the wrong quantity. Second, secure a legitimate
pretrained initialization before tuning anything else; under a tight budget
it is the largest effect measured here. Third, calibrate each client to its
own population (on eICU a single worst-case multiplier would impose $29.34$
on every hospital against a median need of $13.76$), and consider a minimum
cohort size, since the smallest hospitals set the cost. Fourth, use
compression and the edge tier, which cut traffic and energy at no measurable
cost in accuracy, and establish fusion per task. Before deployment, the
end-to-end conditions should hold, with disjoint pretrained encoders and
standardization constants fitted on a disjoint pool, as in the eICU end-to-end
configuration, and the simulated network and modeled energy figures should be
confirmed on physical hardware.

\subsection{Limitations}
\emph{Utility.} On eICU the cost of the guarantee is inconclusive against the
two-point target and exceeds it end to end (Section~\ref{sec:results:privacy});
on WEAR the record-level accuracy loss is inconclusive. \emph{Encoders and
unit.} VitalDB and WEAR utilities use oracle encoders unless stated and are
upper bounds; the disjoint encoders cost $1.57$ AUROC points and leave
differences between configurations unaffected. VitalDB analyses other than the
privacy cost, the cost of the unit, the baseline comparison and the audit are
case-level, covering persons with several operations only by group privacy;
WEAR is record-level; the eICU source partition is qualified in
Section~\ref{sec:method:analysis}; and the validation cohort selects
checkpoints without noise. \emph{Power.} The scaling grids and the recipe
sweep use three seeds, and resolving the eICU gaps of
Table~\ref{tab:kscaling} would need about twenty-five. The audit (sixteen
shadows, VitalDB only) cannot resolve a $0.1\,\%$ false-positive rate; it
bounds these attacks and does not establish the absence of leakage.
\emph{Baselines and simulation.} HED-FL and FL2DP are re-implemented from their
papers, with FL2DP and the baseline audits on VitalDB only, and round
latencies compare simulated topologies rather than a deployed system.
\emph{Fusion and theory.} A fusion gain is shown on VitalDB only, and the
$A(1-1/k)$ form is supported empirically, not proved.
